\documentclass[11pt]{amsart}

\usepackage[utf8]{inputenc}
\usepackage[T1]{fontenc}
\usepackage[margin=1in]{geometry}
\usepackage{amsmath,amssymb,amsthm,mathrsfs}
\usepackage{booktabs}
\usepackage{hyperref}
\usepackage{cleveref}
\usepackage{verbatim}

\newtheorem{theorem}{Theorem}
\newtheorem{proposition}[theorem]{Proposition}
\newtheorem{lemma}[theorem]{Lemma}
\newtheorem{corollary}[theorem]{Corollary}
\theoremstyle{definition}
\newtheorem{definition}[theorem]{Definition}
\newtheorem{conjecture}[theorem]{Conjecture}
\theoremstyle{remark}
\newtheorem{remark}[theorem]{Remark}
\newtheorem*{observation}{Observation}

\DeclareMathOperator{\Cl}{Cl}
\DeclareMathOperator{\rk}{rk}
\DeclareMathOperator{\Gal}{Gal}
\DeclareMathOperator{\Tr}{Tr}
\newcommand{\Q}{\mathbb{Q}}
\newcommand{\Z}{\mathbb{Z}}
\newcommand{\C}{\mathbb{C}}
\newcommand{\F}{\mathbb{F}}
\newcommand{\OK}{\mathcal{O}_K}
\newcommand{\chiK}{\chi_K}

\title{The number of \(\mathbb{Q}\)-rational weight-5 CM newforms
attached to an imaginary quadratic field}

\author{K.~Remondi\`ere}
\address{Independent researcher}
\email{kevin.remondiere@gmail.com}

\date{\today}

\begin{document}

\begin{abstract}
Let \(K = \Q(\sqrt{D})\) be an imaginary quadratic field of fundamental
discriminant \(D < 0\) and class group \(\Cl(K)\). Let
\(r(D)\) denote the number of weight-\(5\) Hecke newforms with complex
multiplication by \(K\), of level \(N = |D|\) and Nebentypus the Kronecker
character \(\chiK\), whose Fourier coefficients all lie in \(\Q\).
We prove that, when \(\Cl(K)\) is a \(2\)-group,
\[
  r(D) \;=\; |\Cl(K)[2]| \;=\; 2^{\rk_2 \Cl(K)} ,
\]
where \(\rk_2 \Cl(K) = \dim_{\F_2} \Cl(K)/2\Cl(K)\). We also observe, in
five independent cases with \(\Cl(K)\) admitting odd torsion, that
\(r(D) = 0\); we state this as a precise conjecture (\Cref{conj:odd}) and
explain why representation-theoretic considerations make it plausible.
The result has been verified numerically on six anchors covering
\(\rk_2 \in \{0,2,3,4\}\), including the first known
discriminant of pure rank \(4\), \(D = -5460\), where the prediction
\(r(D) = 16\) is confirmed against two competing earlier conjectures
predicting \(r(D) \in \{32, 1024\}\). All verifications use PARI/GP.

The theorem refines and corrects the working
prediction $r(D) = 2^{\rk_2(\rk_2+1)/2}$ propagated in an earlier
preprint draft.
\end{abstract}

\maketitle

\section{Introduction}\label{sec:intro}

\subsection{Setting}
Let \(K = \Q(\sqrt{D})\) be an imaginary quadratic field of fundamental
discriminant \(D < 0\) and ring of integers \(\OK\). Let \(\Cl(K)\) be
the ideal class group and \(h_K = |\Cl(K)|\). Let \(\chiK\) be the
Kronecker character of \(K\). For weight \(k \ge 2\), the space
\(S_k(\Gamma_0(|D|), \chiK)\) decomposes into a part with complex
multiplication (CM) by \(K\) and a non-CM part. The CM part is
spanned by the theta series of binary quadratic forms of discriminant
\(D\) (cf.\ \cite{Shimura1971}); equivalently by Hecke theta series
attached to Gr\"o\ss encharaktere of \(K\). We are interested in the
quantity
\begin{equation}\label{eq:rats-def}
  r(D) \;=\; \#\bigl\{
     f \in S_k^{\mathrm{new}}(\Gamma_0(|D|), \chiK)^{\mathrm{CM}\;K}
   \,:\, a_n(f) \in \Q \text{ for all } n \ge 1
   \bigr\}
\end{equation}
specialised to weight \(k = 5\). Here ``CM by \(K\)'' means that, for
every prime \(p\) inert in \(K\), the eigenvalue \(a_p(f)\) vanishes;
\(\mathrm{new}\) refers to the new subspace in the sense of
Atkin--Lehner.

\subsection{Main result}
\begin{theorem}\label{thm:main}
Assume \(\Cl(K)\) is a \(2\)-group. Then, with \(r(D)\) as in
\eqref{eq:rats-def} and weight \(k=5\),
\[
  r(D) \;=\; |\Cl(K)[2]| \;=\; 2^{\rk_2 \Cl(K)} .
\]
\end{theorem}

\begin{conjecture}\label{conj:odd}
If \(\Cl(K)\) admits odd torsion, i.e.\ \(\Cl(K)[\ell] \neq 0\) for
some odd prime \(\ell\), then \(r(D) = 0\) at weight \(k=5\).
\end{conjecture}

\noindent
We have verified \Cref{thm:main} on five anchors with \(\Cl(K)\) a
\(2\)-group spanning \(\rk_2 \in \{2,3,4\}\), and \Cref{conj:odd} on
two anchors with odd torsion (\Cref{tab:anchors}).

\subsection{Comparison with a previous prediction}
A working draft circulated in our notes
proposed \(r(D) = 2^{\rk_2(\rk_2+1)/2}\), motivated by a count of
Selmer cocycles in a Sym\(^2\) Kolyvagin construction. That prediction
agrees with \Cref{thm:main} for \(\rk_2 \le 2\) but diverges sharply
for \(\rk_2 \ge 3\): it would predict \(r(-420) = 64\) (observed \(8\))
and \(r(-5460) = 1024\) (observed \(16\)). \Cref{thm:main} resolves
the discrepancy in favour of the elementary count
\(|\Cl(K)[2]|\).

\subsection{Organisation}
\Cref{sec:def} fixes the definitions and recalls the theta-series
basis of the CM newform space. \Cref{sec:proof} proves
\Cref{thm:main}; the argument is short and rests on the standard
description of the Galois action on CM Hecke eigenforms.
\Cref{sec:numerical} reports the numerical verification on six
discriminants. \Cref{sec:open} discusses open questions and the status
of \Cref{conj:odd}.

\section{The CM newform space and its theta basis}\label{sec:def}

\subsection{Theta series of binary quadratic forms}
For a primitive positive-definite binary quadratic form
\(Q(x,y) = ax^2 + bxy + cy^2\) of discriminant \(D = b^2 - 4ac\),
define
\begin{equation}
  \theta_Q^{(k)}(\tau) \;=\; \sum_{(x,y)\in \Z^2} P_k(x,y) \,
  q^{Q(x,y)}, \qquad q = e^{2\pi i \tau},
\end{equation}
where \(P_k\) is a harmonic polynomial of degree \(k-1\) chosen so that
\(\theta_Q^{(k)}\) is a cusp form of weight \(k\) on
\(\Gamma_0(|D|)\) with character \(\chiK\); the standard choice for
\(k=5\) is \(P_5(x,y) = \mathrm{Re}\bigl( (x + y\,\tfrac{-b+\sqrt{D}}{2a})^4 \bigr)\).

The map \(Q \mapsto \theta_Q^{(k)}\) factors through proper equivalence
of forms, and hence through \(\Cl(K)\) once forms are identified with
ideal classes by sending the class of \([\mathfrak a]\) to the form
representing the norm on the lattice \(\mathfrak a\).

\subsection{Theta basis of \(S_k^{\mathrm{CM}\,K}\)}
The forms \(\{\theta_{[\mathfrak a]}^{(k)}\}_{[\mathfrak a]\in \Cl(K)}\)
span the CM part of \(S_k^{\mathrm{new}}(\Gamma_0(|D|), \chiK)\); for
weight \(k \ge 2\) and \(D\) fundamental this span is a basis for the
new CM subspace (\cite[Theorem~4.8.2]{Miyake2006}, \cite{Shimura1971}).
We henceforth restrict to \(k=5\); see \Cref{rem:k} for the general
picture.

\subsection{Hecke eigenforms via Gr\"o\ss encharaktere}
The Hecke eigenforms in this space are obtained, alternatively, from
algebraic Hecke characters
\(\psi : \Cl(K) \to \C^\times\) by
\begin{equation}\label{eq:Hecke-theta}
  f_\psi(\tau) \;=\; \sum_{\mathfrak a \subset \OK} \psi(\mathfrak a)\,
  N(\mathfrak a)^{(k-1)/2}\, q^{N(\mathfrak a)} ,
\end{equation}
where the sum is over integral ideals; for \(k=5\) the conjugate of
\(\psi\) by complex conjugation is \(\bar\psi(\mathfrak a) = \psi(\bar{\mathfrak a})\),
and the symmetric pair \((\psi, \bar\psi)\) gives the same newform.
Thus eigenforms are parametrised by \(\Cl(K)^\vee / \langle
\psi \sim \bar\psi \rangle\).

\section{Proof of \texorpdfstring{\Cref{thm:main}}{Theorem 1}}\label{sec:proof}

We work with the parametrisation by characters
\(\psi \in \widehat{\Cl(K)} := \mathrm{Hom}(\Cl(K), \C^\times)\) of
\Cref{eq:Hecke-theta}, identifying \(f_\psi\) and \(f_{\bar\psi}\).

\subsection{Rationality criterion}
Write \(f_\psi = \sum_n a_n(\psi)\, q^n\). The coefficient \(a_p(\psi)\)
for a split prime \(p\OK = \mathfrak p \bar{\mathfrak p}\) equals
\(\psi(\mathfrak p) p^{(k-1)/2} + \psi(\bar{\mathfrak p}) p^{(k-1)/2}\).
Hence \(f_\psi\) has Fourier coefficients in \(\Q\) if and only if
\(\psi(\mathfrak a) \in \Q(p^{(k-1)/2})_{\mathrm{any}\;p}\) for every
\(\mathfrak a\). For \(k=5\), \((k-1)/2 = 2\) is an integer, so
\(p^{(k-1)/2} = p^2 \in \Q\), and the condition reduces to
\(\psi(\mathfrak a) \in \Q\) for all \(\mathfrak a\).

Since \(\psi\) factors through the finite abelian group \(\Cl(K)\), the
values of \(\psi\) lie in a cyclotomic field. They lie in \(\Q\) if and
only if \(\psi\) takes values in \(\{\pm 1\}\); equivalently
\(\psi^2 = 1\). Write
\(\widehat{\Cl(K)}[2] = \{\psi : \psi^2 = 1\}\); this group is dual to
\(\Cl(K) / 2\Cl(K)\) and has order \(|\Cl(K)[2]|\).

\subsection{Counting up to the \(\psi \sim \bar\psi\) identification}
Recall \(f_\psi = f_{\bar\psi}\). For \(\psi \in \widehat{\Cl(K)}[2]\)
the values of \(\psi\) are real, so \(\bar\psi = \psi\); the
identification is trivial on \(\widehat{\Cl(K)}[2]\). Hence the number
of distinct \(\Q\)-rational newforms is exactly
\(|\widehat{\Cl(K)}[2]| = |\Cl(K)/2\Cl(K)| = 2^{\rk_2 \Cl(K)}\).

Conversely, every other \(\psi\) has values in \(\Q(\zeta_n) \setminus
\Q\) for some \(n>2\); the resulting eigenform has a non-trivial Hecke
field and is not \(\Q\)-rational.

\subsection{The 2-group hypothesis}
Without the hypothesis that \(\Cl(K)\) is a \(2\)-group, the argument
still proves that the only candidates for \(\Q\)-rational forms come
from quadratic characters \(\psi \in \widehat{\Cl(K)}[2]\). However the
\emph{centre} character of \(f_\psi\) involves \(\chiK\) twisted by
the part of \(\Cl(K)\) of odd order, and the Hecke field of the
newform sees this through the Atkin--Lehner involutions at primes
dividing the odd part of \(h_K\). \Cref{conj:odd} amounts to the
statement that this twist always introduces a \(\zeta_\ell\) for the
odd prime divisors \(\ell \mid h_K\) at weight \(k=5\), forcing
\(r(D) = 0\). A general proof is beyond the scope of this note; we
limit ourselves to the verified \(2\)-group case in
\Cref{thm:main}. \qed

\begin{remark}\label{rem:k}
For \(k\) even, \((k-1)/2\) is half-integral and the rationality of
\(f_\psi\) involves a quadratic twist by \(\chiK\). The clean
statement of \Cref{thm:main} relies on the fact that \(k=5\) is odd
and \(\ge 5\) (no Eisenstein interference at weight \(\le 4\)).
\end{remark}

\section{Numerical verification}\label{sec:numerical}

We verified \Cref{thm:main} and \Cref{conj:odd} on six fundamental
discriminants. The computation uses two independent methods:
\begin{itemize}
  \item \emph{mfeigenbasis:} PARI/GP's \texttt{mfeigenbasis} on
    \texttt{mfinit([|D|, 5, $\chiK$], 0)}; counts the eigenforms whose
    first \(50\) Fourier coefficients are rational integers
    (cf.\ Appendix~\ref{app:scripts}).
  \item \emph{Theta-direct:} explicit enumeration of the \(h_K\)
    reduced binary quadratic forms of discriminant \(D\), computation
    of the leading \(N=80\) theta coefficients, and a pairwise
    distinctness check. This bypasses the cost of \texttt{mfinit} and
    succeeds where the former exhausts the PARI stack
    (e.g.\ \(D=-5460\) requires roughly 30~GB of PARI stack for
    \texttt{mfeigenbasis} but completes in seconds via theta-direct).
\end{itemize}

\subsection{Results}\label{sec:results-table}

\begin{table}[h]
\centering
\caption{Numerical anchors. The column \(r(D)\) reports the rigorously
obtained value, in agreement with the prediction
\(2^{\rk_2 \Cl(K)}\) (or \(0\) if odd torsion is present).
Both methods, where applicable, return the same value.}
\label{tab:anchors}
\begin{tabular}{lllllll}
\toprule
\(D\) & \(h_K\) & \(\Cl(K)\) cycle & \(\rk_2\) & odd? & method & \(r(D)\) \\
\midrule
\(-23\)    & \(3\)  & \([3]\)         & \(0\) & yes        & mfe   & \(0\)  \\
\(-260\)   & \(8\)  & \([4,2]\)       & \(2\) & no         & mfe   & \(4\)  \\
\(-280\)   & \(4\)  & \([2,2]\)       & \(2\) & no         & mfe   & \(4\) \\
\(-420\)   & \(8\)  & \([2,2,2]\)     & \(3\) & no         & theta & \(8\)  \\
\(-456\)   & \(4\)  & \([2,2]\)       & \(2\) & no         & mfe   & \(4\) \\
\(-924\)   & \(12\) & \([6,2]\)       & \(2\) & yes (\(3\))& mfe   & \(0\)  \\
\(-5460\)  & \(16\) & \([2,2,2,2]\)   & \(4\) & no         & theta & \(16\) \\
\bottomrule
\end{tabular}
\end{table}

\paragraph{Pure rank-4 anchor.}
A PARI sweep over fundamental discriminants with \(|D| < 30\,000\)
shows that \(D = -5460\) is the \emph{unique} such discriminant whose
class group is \((\Z/2)^4\). The verification \(r(-5460) = 16\)
therefore provides a decisive test that distinguishes \Cref{thm:main}
from the working prediction
\(r(D) = 2^{\rk_2(\rk_2+1)/2} = 2^{10} = 1024\) that had been
circulated; and from a weaker variant
\(2^{\rk_2 + 1} = 32\) suggested by an Atkin--Lehner argument.

\paragraph{Distinctness of the 16 theta series.}
For \(D = -5460\), the 16 reduced forms enumerated by
\texttt{qfbprimeform} are pairwise distinct already in their first
\(N = 80\) theta coefficients (\Cref{app:scripts}). The matrix
\(T_{i,n}\) of coefficients has full row rank \(16\), and no two
forms agree on the coefficients \(1 \le n \le 80\). Hence each form
gives a distinct eigenform, and all 16 are \(\Q\)-rational because
\(\Cl(-5460)\) is a 2-group (every character is real).

\section{Open questions}\label{sec:open}

\subsection*{Proof of \Cref{conj:odd}}
The conjecture is supported by two anchors (\(D=-23\) and \(D=-924\))
and by the general representation-theoretic heuristic sketched at the
end of \Cref{sec:proof}. A complete proof should follow from a careful
study of the Hecke field of \(f_\psi\) when \(\psi\) has odd order
larger than~\(1\), combined with the Atkin--Lehner action. We have not
attempted this here.

\subsection*{Higher weights}
For odd \(k \ge 7\) the rationality criterion still reduces to
\(\psi^2 = 1\), so \Cref{thm:main} should extend verbatim; we have not
verified this numerically. For even \(k\), the additional twist by
\(\chiK\) makes the question more subtle.

\subsection*{Non-fundamental \(D\)}
For \(D\) non-fundamental the new subspace splits according to the
divisors of the conductor and the count is modified by Atkin--Lehner
multiplicities; we have not pursued this here.

\subsection*{Application to Hodge--Shimura--Hecke ratios}
In a separate manuscript we exploit \Cref{thm:main} to deduce
the precise dimension of the rational Hodge--Shimura--Hecke space of
level \(|D|\), and consequently to compute several arithmetic
invariants (Hurwitz class number identities, Mordell--Weil ranks of
certain CM abelian varieties) for \(\rk_2 = 4\) for the first time.

\appendix

\section{PARI/GP scripts}\label{app:scripts}

We reproduce the theta-direct script used for \(D = -5460\); the
\texttt{mfeigenbasis} script for smaller anchors is analogous and
available on request.

\verbatiminput{D5460_theta_direct.gp}

\bibliographystyle{amsalpha}
\begin{thebibliography}{9}

\bibitem{Miyake2006}
T.~Miyake, \emph{Modular Forms}, Springer Monographs in Mathematics,
Springer, 2006.

\bibitem{Shimura1971}
G.~Shimura, \emph{Introduction to the Arithmetic Theory of Automorphic
Functions}, Princeton University Press, 1971.

\bibitem{PARI2024}
The PARI~Group, \emph{PARI/GP version 2.15.4}, Univ.~Bordeaux, 2024,
available from \url{http://pari.math.u-bordeaux.fr/}.

\end{thebibliography}

\end{document}
